\documentclass[aspectratio=169]{beamer}

% Match the existing CSE 534 Beamer styling used in the N-gram deck
\usetheme{default}
\usecolortheme{default}
\setbeamertemplate{navigation symbols}{}
\setbeamertemplate{footline}[frame number]

\definecolor{darkbg}{HTML}{000000}
\definecolor{lighttext}{HTML}{999999}
\definecolor{miamired}{HTML}{941728}
\definecolor{codebg}{HTML}{1a1a1a}

\setbeamercolor{background canvas}{bg=darkbg}
\setbeamercolor{normal text}{fg=lighttext}
\setbeamercolor{frametitle}{fg=miamired}
\setbeamercolor{title}{fg=miamired}
\setbeamercolor{structure}{fg=miamired}
\setbeamercolor{item}{fg=miamired}
\setbeamercolor{block title}{fg=miamired,bg=codebg}
\setbeamercolor{block body}{fg=lighttext,bg=codebg}

\setbeamerfont{title}{series=\bfseries,size=\Large}
\setbeamerfont{frametitle}{series=\bfseries}

\usepackage{graphicx}
\usepackage{listings}
\usepackage{tikz}
\usepackage{array}
\usepackage{booktabs}
\usepackage{colortbl}
\usepackage{amsmath}
\usepackage{amssymb}

\lstset{
    basicstyle=\ttfamily\small\color{lighttext},
    backgroundcolor=\color{codebg},
    keywordstyle=\color[HTML]{569cd6},
    commentstyle=\color[HTML]{6a9955},
    stringstyle=\color[HTML]{ce9178},
    showstringspaces=false,
    breaklines=true,
    frame=single,
    rulecolor=\color{codebg}
}

\setlength{\arrayrulewidth}{0.5pt}
\arrayrulecolor{miamired}
\renewcommand{\textbf}[1]{\textcolor{white}{\bfseries #1}}

\title{Likelihood, Log-Likelihood, and Maximum Likelihood}
\subtitle{Why machine learning uses \emph{log-likelihood}}
\author{CSE 534}
\date{}

\begin{document}

\begin{frame}
\titlepage
\begin{center}
\textbf{Goal:} compare the probability of data with the plausibility of model parameters
\end{center}
\end{frame}

\begin{frame}[c]{Probability vs. likelihood}
\textbf{Probability:} parameters are known; ask what data may appear.

\vspace{0.5em}
\textbf{Likelihood:} data are known; ask which parameters make the data most plausible.

\vspace{1em}
\textbf{Example: coin with } $\mu = P(\text{heads})$
\begin{itemize}
    \item \textbf{Probability:} if $\mu=0.70$, what fraction of heads might we see?
    \item \textbf{Likelihood:} after seeing 70 heads in 100 flips, which $\mu$ is most plausible?
\end{itemize}
\end{frame}

\begin{frame}[c]{Independent observations}
For i.i.d. observations $x_1, \ldots, x_n$:
\[
L(\theta; x_1, \ldots, x_n) = \prod_{i=1}^n p(x_i \mid \theta)
\]

\vspace{0.5em}
\textbf{Bernoulli example:}
\[
L(\mu) = \mu^k (1-\mu)^{n-k}
\]
where $k$ is the number of 1s and $n$ is the total number of observations.
\end{frame}

\begin{frame}[c]{The underflow problem}
\textbf{Goal:} maximize likelihood.

\vspace{0.5em}
But multiplying many probabilities drives the product toward zero:
\[
L(\mu=0.7)=0.7^{70000}\cdot 0.3^{30000} \rightarrow 0
\]

\vspace{0.5em}
\textbf{Outcome:} all candidates underflow to zero, so we cannot compare them.

\vspace{0.5em}
\textbf{Demo A:} \texttt{05a\_likelihood\_underflow.py}
\end{frame}

\begin{frame}[c]{Log-likelihood solves it}
\textbf{Key idea:} products become sums under a logarithm.
\[
\ell(\theta)=\log L(\theta)=\sum_{i=1}^n \log p(x_i\mid \theta)
\]

\vspace{0.5em}
\textbf{Why this helps:}
\begin{itemize}
    \item maximizing $L$ is equivalent to maximizing $\log L$
    \item addition is numerically stable
    \item gradients remain usable for optimization
\end{itemize}

\vspace{0.5em}
For Bernoulli data:
\[
\ell(\mu)=\sum_{i=1}^n \left[x_i\log\mu + (1-x_i)\log(1-\mu)\right]
\]
\end{frame}

\begin{frame}[c]{Maximum likelihood estimate}
We want the value of $\mu$ that maximizes $\ell(\mu)$.

\vspace{0.5em}
For Bernoulli observations, the maximum-likelihood estimate is the sample mean:
\[
\hat{\mu}_{ML} = \frac{1}{n}\sum_{i=1}^n x_i
\]

\vspace{0.5em}
This is exactly the observed fraction of 1s.

\vspace{0.5em}
For 70 heads in 100 flips:
\[
\hat{\mu}_{ML} = 70/100 = 0.70
\]
\end{frame}

\begin{frame}[c]{The derivation}
Start from:
\[
\ell(\mu)=\sum_{i=1}^n \left[x_i\log\mu + (1-x_i)\log(1-\mu)\right]
\]

Differentiate with respect to $\mu$:
\[
\frac{\partial \ell}{\partial \mu}
=\sum_{i=1}^n\left[\frac{x_i}{\mu} - \frac{1-x_i}{1-\mu}\right]
\]

Set it to zero and solve:
\[
\hat{\mu}_{ML}=\frac{1}{n}\sum_{i=1}^n x_i
\]
\end{frame}

\begin{frame}[c]{Gradient vanishing}
Modern ML uses gradient-based optimization.
\[
\theta \leftarrow \theta + \alpha \cdot \nabla L(\theta)
\]

\vspace{0.5em}
When raw likelihood becomes extremely small, the gradient collapses:
\begin{itemize}
    \item $\nabla L(\theta) \rightarrow 0$
    \item parameters stop updating
    \item optimization fails
\end{itemize}

\vspace{0.5em}
\textbf{Demo B:} \texttt{05b\_likelihood\_gradient\_fail.py}
\end{frame}

\begin{frame}[c]{Log-likelihood succeeds}
The same idea works if we optimize the log-likelihood instead:
\[
\frac{\partial \ell}{\partial \mu}
=\sum_{i=1}^n\left[\frac{x_i}{\mu} - \frac{1-x_i}{1-\mu}\right]
\]

\vspace{0.5em}
\textbf{Demo C:} \texttt{05c\_likelihood\_logit.py}

\vspace{0.5em}
This converges to $\mu \approx 0.70$.
\end{frame}

\begin{frame}[c]{Negative log-likelihood}
Machine learning usually minimizes a loss function instead of maximizing a likelihood.

\[
\text{Loss} = -\ell(\theta) = -\log L(\theta)
\]

\vspace{0.5em}
\textbf{Key idea:}
\begin{itemize}
    \item maximizing likelihood $\equiv$ minimizing negative log-likelihood
    \item classification uses cross-entropy loss
    \item regression often uses Gaussian NLL / MSE equivalence
\end{itemize}
\end{frame}

\begin{frame}[c]{Why this matters for deep learning}
Large language models are trained by minimizing the negative log-likelihood of the observed text:
\[
\text{minimize } -\log P(\text{text} \mid \theta)
\]

\vspace{0.5em}
This is the reason the log term appears everywhere in modern ML:
\begin{itemize}
    \item more stable numerical behavior
    \item better gradients
    \item direct connection to model likelihood
\end{itemize}
\end{frame}

\begin{frame}[c]{Takeaway}
\textbf{Likelihood} asks: given parameters, how likely was the observed data?

\textbf{Log-likelihood} asks: which parameters make the observed data most likely, without underflow?

\vspace{0.5em}
\textbf{Maximum-likelihood estimation} chooses the parameter values that maximize the observed data probability.

\vspace{0.5em}
\textbf{Deep learning} trains by minimizing negative log-likelihood, which is the same optimization problem in a numerically stable form.
\end{frame}

\end{document}
